\documentclass[../../main.tex]{subfiles}
\begin{document}
\begin{oframed}
  \subfile{problem}
\end{oframed}

{\bf [解答]}

\paragraph{(1)}

与えられた方程式
\begin{align}
  \int_0^x{(f(t))}^{-a}dt=\int_a^{f(x)}(e^{\displaystyle-\frac{t^2}{2}}+t^{-a})dt \label{eq:1}
\end{align}
の両辺を $x$ で微分し，$y=f(x)$ を代入して
\begin{align}
  y^{-a} = \left( e^{-\frac{y^2}{2}} + y^{-a} \right) \cdot \dfrac{dy}{dx} \label{eq:2}
\end{align}
である．

次に\cref{eq:1}に $x=0$ を代入して
\begin{align}
  0 = \int_a^{f(0)} \left( e^{-\frac{t^2}{2}} + t^{-a} \right) dt \label{eq:3}
\end{align}
である．右辺の被積分関数は $t>0$ で正だからこれが成り立つのは $f(0) = a$ の時である($\because f(x)>0$)．

\paragraph{(2)-(イ)}

\cref{eq:2}から
\begin{align}
  y' = \dfrac{1}{y^a e^{-\frac{y^2}{2}} + 1} \quad (\because y>0) \label{eq:4}
\end{align}
である．$g(y) = y^a e^{-\frac{y^2}{2}}$ とおくと
\begin{align}
  y' = \dfrac{1}{g(y)+1}
\end{align}
であって，一階微分は
\begin{align}
  g'(y) = e^{-\frac{y^2}{2}} y^{a-1} (a - y^2)
\end{align}
となるから，$g(y)$の増減表は以下のようになる．

\begin{table}[htb]
  \centering
  \caption{$g(y)$ の増減表}
  \begin{tabular}{c|c|c|c|c|c}
    $y$ & 0 & $\cdots$ & $\sqrt{a}$ & $\cdots$ & \\
    \hline
    $g'$ & & $+$ & $0$ & $-$ & \\
    \hline
    $g$ & 0 & $\nearrow$ & $a^{\frac{a}{2}} e^{-\frac{a}{2}}$ & $\searrow$ &
  \end{tabular}
\end{table}

従って\cref{eq:4}において$0 < y$ の時
\begin{align}
  0 < \dfrac{1}{a^{\frac{a}{2}} e^{-\frac{a}{2}} + 1} \le \dfrac{1}{g(y)+1} = y' \le 1
\end{align}
だから $b = \dfrac{1}{a^{\frac{a}{2}} e^{-\frac{a}{2}} + 1}$ とおけば ($b>0$) $b \le y' \le 1$ なる $b$ が存在する．\qed

\paragraph{(2)-(ロ)}

\cref{eq:4}から
\begin{align}
  y \left(\dfrac{1}{y'} - 1\right) = y^{a+1} e^{-\frac{y^2}{2}} \label{eq:5}
\end{align}
となる．この右辺は$g(y)$ において $a$ を $a+1$ に置きかえたものだから，(イ)より$y>0$ の時
\begin{align}
  0 < y \left(\dfrac{1}{y'} - 1\right) \le (a+1)^{\frac{a+1}{2}} e^{-\frac{a+1}{2}}
\end{align}
となる．
よって$c = (a+1)^{\frac{a+1}{2}} e^{-\frac{a+1}{2}}$ とおけば，$0 \le y \left(\dfrac{1}{y'} - 1\right) \le c$ なる $c$ が存在する．\qed

\paragraph{(3)}

(2)(イ) の不等式$b \le f'(x) \le 1$の両辺を$[0,x]$で積分して
\begin{align}
  & \left[ bt \right]_0^x \le \left[ f(t) \right]_0^x \le \left[ t \right]_0^x \\
  & bx \le f(x) - f(0) \le x \\
  & bx + a \le f(x) \le x + a \quad (\because f(0)=a) \label{eq:6}
\end{align}
である．$b>0$ からはさみうちの定理を用いて
\begin{align}
  \lim_{x \to \infty} f(x) = \infty \label{eq:7}
\end{align}
を得る．

そこで\cref{eq:4}より求める極限値は
\begin{align}
  \lim_{x \to \infty} f'(x) = \lim_{y \to \infty} \dfrac{1}{y^a e^{-\frac{y^2}{2}} + 1} = 1
\end{align}
である．

次に$f'(x)$の最小値を求める．\cref{eq:4}より$f(x)$の挙動がわかれば$f'(x)$の挙動がわかるのでこれを調べる．(イ)から$f'(x)>0$だから$f(x)$ は単調増加である．これと $f(0) = a$ から，$a \le f(x)$ である．
これと $\lim_{x\to\infty} f(x) = \infty$ に注意して，(イ)の増減表から以下の二つの場合が考えられる．

\subparagraph{$1^\circ \ a \le \sqrt{a} \Leftrightarrow 0 < a \le 1$ の時}

\begin{table}[htb]
  \centering
  \caption{$0 < a \le 1$ における $g(y)$ の増減表 ($y \ge a$)}
  \label{tab:variation_case1}
  \begin{tabular}{c|ccccc}
    $y$ & $a$ & $\cdots$ & $\sqrt{a}$ & $\cdots$ \\
    \hline
    $g'$ & & $+$ & $0$ & $-$ \\
    \hline
    $g$ & $a^a e^{-\frac{a^2}{2}}$ & $\nearrow$ & $a^{\frac{a}{2}} e^{-\frac{a}{2}}$ & $\searrow$
  \end{tabular}
\end{table}

この時は増減表より$g(y)$は$y=\sqrt{a}$で最大だから，\cref{eq:4}よりこの時$f'(x)$が最小で
\begin{align}
  \min y' = \dfrac{1}{g(\sqrt{a}) + 1} = \dfrac{1}{a^{\frac{a}{2}} e^{-\frac{a}{2}} + 1} \label{eq:8}
\end{align}
である．

\subparagraph{$2^\circ \ \sqrt{a} \le a \Leftrightarrow 1 \le a$ の時}

\begin{table}[htb]
  \centering
  \caption{$1 \le a$ における $g(y)$ の増減表 ($y \ge a$)}
  \label{tab:variation_case2}
  \begin{tabular}{c|ccc}
    $y$ & $a$ & $\cdots$ \\
    \hline
    $g'$ & & $-$ \\
    \hline
    $g$ & $a^a e^{-\frac{a^2}{2}}$ & $\searrow$
  \end{tabular}
\end{table}

この時は増減表より端点$y=a$で$g(y)$が最大だから，\cref{eq:4}よりこの時$f'(x)$が最小で
\begin{align}
  \min y' = \dfrac{1}{g(a) + 1} = \dfrac{1}{a^a e^{-\frac{a^2}{2}} + 1} \label{eq:9}
\end{align}
である．
\casesend

以上\cref{eq:8,eq:9}をまとめて求める最小値は
\begin{align}
  \begin{cases}
    0 < a \le 1 \text{ の時} & \dfrac{e^{\frac{a}{2}}}{a^{\frac{a}{2}} + e^{\frac{a}{2}}} \\
    1 \le a \text{ の時} & \dfrac{e^{\frac{a^2}{2}}}{a^a + e^{\frac{a^2}{2}}}
  \end{cases}
\end{align}
である．

{\bf [別解]}

\paragraph{(3)}

(ロ)を直接使っても良い．(ロ)の両辺$f(x)$で割って
\begin{align}
  0 \le \dfrac{1}{f'(x)} - 1 \le \dfrac{c}{f(x)}
\end{align}
だから，$x\to\infty$で$f(x)\to\infty$であることから挟み撃ちの定理から
\begin{align}
  \lim_{x\to\infty} f'(x) = 1
\end{align}
である．

\end{document}
